\section{A differential inequality for the disc envelope}
\label{sec:variation}

Let $v$ be the lower semicontinuous envelope defined in
\eqref{var:envelope}--\eqref{var:sandwich}, with its fixed horizon $T$.
We continue to denote physical time at a disc center by $s$.
The following inequality is the input to the fiber optimization in
Section~\ref{sec:ellipsoids}.

\begin{proposition}[Differential inequality for the envelope]
\label{prop:disc-inequality}
Suppose that a smooth real function $\varphi$ touches $v$ from below at
$(y_*,s_*)$, where $0<s_*<T$.  In holomorphic cotangent coordinates write
\[
 W=\dd_Y\varphi(y_*,s_*)=
 \begin{pmatrix} A&B\\ B^*&C\end{pmatrix},\qquad C>0.
\]
There is a Hermitian form $K$ on the horizontal coordinate space such that
\begin{equation}
 K\ge0,\qquad K\ge A-BC^{-1}B^*,\qquad
 \partial_s\varphi(y_*,s_*)\ge\tr(C^{\mathsf T}K).
 \label{var:viscosity-inequality}
\end{equation}
The transpose is the tangent--cotangent contraction in these coordinates.
In particular the statement applies to lower tests translated in a fixed
cotangent coordinate chart; its contraction has constant coefficients.
\end{proposition}

We prove the Proposition through localized almost minimizing discs.
The estimate on their deformations will be stated with its constants:
the compact set containing the interiors of the discs is allowed to
depend on the area penalty, whereas the coefficient multiplying that
penalty will depend only on boundary data.
After this deformation estimate, we remove the negative part of the
boundary Hessian and factor the remaining positive form.  A backward
clock variation bounds its symplectic cost.  The final Schur-complement
limit gives the asserted differential inequality.

\subsection{Localization and the order of approximation}

Subtracting a positive multiple of
$|y-y_*|^4+|s-s_*|^4$ from $\varphi$ makes the contact strict on a compact
coordinate cylinder $\mathcal C\Subset Y\times(0,T)$ without changing
the derivatives in \eqref{var:viscosity-inequality}.
Thus $v-\varphi\ge0$ on $\mathcal C$, with equality only
at $(y_*,s_*)$, and with a positive gap near its boundary.  The definition
of lower semicontinuous regularization gives
\[
 \inf_{(y,s)\in\mathcal C}\bigl(v^{\rm raw}(y,s)-\varphi(y,s)\bigr)=0.
\]

On an interval slightly larger than $[0,T]$, take $Q$ and $\rho_j$ from
Theorem~\ref{thm:localized-flow} and
Proposition~\ref{prop:scalar-comparison}, and set
\begin{equation}
 d_{0,j}(x,t)=j^{-1}Q(x,t)+\rho_j(x,t).
 \label{var:base-cost}
\end{equation}
Thus $d_{0,j}\ge0$, $(\partial_t-\Delta_{h_t})d_{0,j}\ge0$, and
$d_{0,j}\to0$ uniformly on compact space--time sets.  Choose increasing
$A_j\to\infty$ so that
\begin{equation}
 d_{0,j}(x,t)\ge\rho(x,t)+j\quad\hbox{if }u(x)>A_j/3.
 \label{var:tail-wall}
\end{equation}
This is possible because $\rho_j\ge\rho-C_j$ and $Q$ has a positive
quadratic lower bound.  Choose fixed smooth scalar profiles and set
$\delta_1=\chi(u/A_j)$ and
$\delta_2=\widetilde\chi(N(\cdot,0)/B_j^2)$, with
$0\le\delta_i\le1$, where $\delta_1>0$ exactly when $u<A_j$ and
$\delta_1=1$ when $u\le2A_j/3$, while $\delta_2>0$ exactly when
$N(y,0)<B_j^2$ and $\delta_2=1$ when $N(y,0)\le B_j^2/2$.
Here $B_j\ge j$ will be increased below.  Smooth cutoffs flat at their
zero sets are permissible.  On the open boundary region put
\begin{equation}
 k_i=\delta_i^{-1}-1,\qquad
 d(y,t,s)=d_{0,j}(x,t)+(1+s)(k_1(x)+k_2(y)).
 \label{var:wall-cost}
\end{equation}
Only the boundary of a competing disc is subject to these spatial
restrictions.  Their closed boundary region is a compact set $K_j\subset Y$.
Each fixed $d_{0,j}$ is smooth one-sided at $t=0$, with all derivatives
bounded on compact sets.  When needed, use a smooth extension matching
these initial jets; extending it by zero to negative time is unnecessary.

Use the boundary seminorm of Section~\ref{sec:discs},
\[
 D(a)^2=\sum_{k\in\mathbb Z}|k|\,\|a_k\|^2,
\]
and add the nonnegative energy
\begin{equation}
 \begin{split}
 \mathcal E(f,w)={}&\int_\Delta\bigl(x_f^*\omega_g+
 f^*G^*\omega_{\rm eucl}+
 (\mathfrak g\circ f)^*\omega_{\Gr}
 +w^*\omega_{\rm eucl}\bigr)\\
 &+D\bigl(\log(\delta_1\delta_2)\circ f\bigr)^2.
 \end{split}
 \label{var:energy}
\end{equation}
Here $x_f=\pi\circ f$, and $\mathfrak g$ is the Gauss map of the immersion $G$ from
Lemma~\ref{lem:disc-frames}.
Let $\mu_{j,\eps}$ be the infimum, over centers in $\mathcal C$ and discs
with the stated boundary restrictions, of
\begin{equation}
 \mathcal F_{j,\eps}(f,w,s)=
 \langle N(f,t)+d(f,t,s)\rangle+\eps\mathcal E(f,w)
 -\varphi(f(0),s),\qquad t=s e^{w+\bar w}.
 \label{var:penalized-functional}
\end{equation}
Infima are sufficient; no minimizing disc is asserted to exist.
Every fixed admissible smooth disc has finite energy.  Also every fixed
disc in \eqref{var:envelope} eventually has $\delta_1=\delta_2=1$ on its
boundary.  Consequently
\begin{equation}
 \mu_{j,0}\longrightarrow0,\qquad
 \mu_{j,\eps}\downarrow\mu_{j,0}\quad(\eps\downarrow0).
 \label{var:infimum-limits}
\end{equation}
For a minimizing sequence indexed by $\nu$, at fixed $j,\eps>0$,
nonnegativity of the unpenalized excess gives
\begin{equation}
 \limsup_{\nu}\eps\mathcal E_\nu
 \le\mu_{j,\eps}-\mu_{j,0},\qquad
 \lim_{\eps\downarrow0}\limsup_\nu\eps(1+\mathcal E_\nu)=0.
 \label{var:vanishing-penalty}
\end{equation}
As $j\to\infty$ after these limits, the centers approach $(y_*,s_*)$ and
\begin{equation}
 \langle d\rangle\longrightarrow0,\qquad
 \langle N(f,t)\rangle\le C.
 \label{var:small-walls}
\end{equation}
Indeed each of $v-\varphi$, the boundary cost $d$, and $\eps\mathcal E$
is nonnegative and bounded by \eqref{var:penalized-functional}.
The strict contact and the boundary gap in $\mathcal C$ imply the
assertion about centers.  They therefore have a uniform interior margin
for large $j$, small $\eps$, and late terms of the sequences.

At fixed $j,\eps$ the energies are bounded.  Proposition~\ref{prop:disc-compactness}
then confines all images of $f$ to a common compact $K_{j,\eps}'$.
There is no claim that $K_{j,\eps}'$ is independent of $\eps$.
One may require the initial discs to extend holomorphically past the
circle: replacing $(f,w)$ by $(f(r\zeta),w(r\zeta))$, $r\uparrow1$,
preserves their centers and converges in every boundary norm used here.
For each disc its strict boundary inequalities persist for $r$ close
enough to one.  Choosing the approximation separately for each term
preserves all infima and minimizing-sequence assertions above.

\subsection{Boundary calculus and holomorphic deformations}

We first give the boundary estimates used to distinguish the two kinds
of constants.  The difference-quotient formula for $D$ gives
\begin{equation}
 \begin{split}
 D(ab)&\le\|a\|_\infty D(b)+\|b\|_\infty D(a),\\
 D(\Psi(a))&\le\Lip(\Psi)D(a),\\
 |\langle a,\partial_\theta b\rangle|&\le D(a)D(b).
 \end{split}
 \label{var:boundary-calculus}
\end{equation}
The last inequality follows directly from Fourier series; it is the
pairing between $\dot H^{1/2}$ and $\dot H^{-1/2}$.  All three formulas
hold componentwise for finite-dimensional vectors and matrices, with
fixed dimensional constants.  Smooth functions of bounded data can
therefore be composed and multiplied with constants determined by their
smooth bounds on the relevant compact sets.  The circular Hilbert
transform preserves $D$ on mean-zero functions.

For a disc let $E_f$ be the frame in Lemma~\ref{lem:disc-frames}, and put
\[
 X=G\circ f,\qquad U=dG\,E_f,\qquad p=(\delta_1\delta_2)\circ f,
 \qquad q=\log p.
\]
In the boundary estimates below, $p,q$ denote these scalar functions.
The symplectic polynomial $p(\mathcal H)$ is distinguished by its matrix
argument.  For an integer
$a\ge4$, to be enlarged once below, let $m=p^a$ and let $h$ be the scalar
outer function with boundary modulus $m$ and $h(0)>0$.  Thus
\begin{equation}
 h(0)=\exp\langle\log m\rangle,
 \quad \|h\|_\infty\le1,
 \quad D(h)+D(h/m)\le C_a D(q).
 \label{var:outer-wall}
\end{equation}
Here $h/m$ denotes only its boundary phase.  To verify the last bound,
write that phase as $\exp(i\mathcal H(aq))$; both the exponential on the
imaginary axis and $q\mapsto e^{aq}$ for $q\le0$ are Lipschitz.
The frame and energy estimates give
\begin{equation}
 \|U\|_\infty\le1,\qquad
 D(X)+D(U)+D(w)+D(q)\le C(1+\sqrt{\mathcal E}).
 \label{var:basic-boundary-bounds}
\end{equation}
Constants absorbing fixed area normalizations are harmless.  Since
$t\le T$, the exponential restricted to the range of $2\Re w$ gives
$D(t)\le C_TD(w)$, even if some boundary times approach zero.

For $t$ and $s$ held fixed in the spatial Hessian, define
\begin{equation}
 W_b=\dd_Y(N+d),\qquad T_b=m^2E_f^*W_bE_f.
 \label{var:boundary-symbol}
\end{equation}
For sufficiently large fixed $a$,
\begin{equation}
 \|T_b\|_\infty\le C_j,
 \qquad D(T_b)^2\le C_j(1+\mathcal E).
 \label{var:symbol-bounds}
\end{equation}
For example
$\dd(\delta^{-1})=2\delta^{-3}\partial\delta\otimes\bar\partial\delta
-\delta^{-2}\dd\delta$.
Multiplication by $p^{2a}$ clears all its denominators and leaves smooth
coefficients on the closed boundary region.  Every other coefficient
is smooth on $K_j\times[0,T]\times\mathcal C_s$.
Choose $K_j\subset\mathcal W_j\Subset Y$ with compact closure.
The injective immersion $G$ restricted to $\overline{\mathcal W_j}$
has a continuous inverse onto its image by compactness; thus
$G|_{\mathcal W_j}$ is an embedding.
A finite extension and partition construction from $G(\mathcal W_j)$
near $G(K_j)$ expresses these coefficients as bounded smooth functions
of $X$ there.
The same applies to tensor coefficients evaluated on $U$.
Equation~\eqref{var:symbol-bounds} now follows from
\eqref{var:boundary-calculus} and \eqref{var:basic-boundary-bounds}.
These coefficient extensions are used only along the embedded compact
set; no properness of $G$ is asserted.

\begin{lemma}[Deformation estimate, with its uniformities]
\label{lem:disc-deformation}
Use the preceding localization, with $j$ sufficiently large and
$\eps>0$ sufficiently small that the minimizing-sequence centers have
a common interior margin in $\mathcal C$.  Fix a bound $L<\infty$
and such a sequence for \eqref{var:penalized-functional}, with energy
bounded by $E_0<\infty$.  Let $l_\nu$ be any holomorphic columns, smooth on the
closed disc, such that
\[
 \|l_\nu\|_\infty\le L,\qquad D(l_\nu)\le D_0<\infty.
\]
They may depend on the disc.  With
$\varphi_{E,c}=E_f(0)^*\dd_Y\varphi(f(0),s)E_f(0)$, one has
\begin{equation}
 \limsup_{\nu\to\infty}
 \left\{
 h(0)^2\varphi_{E,c}(l(0),\overline{l(0)})
 -\langle T_b(l,\bar l)\rangle
 -C_j(L)\eps\bigl(1+\mathcal E+D(l)^2\bigr)
 \right\}\le0.
 \label{var:deformation-estimate}
\end{equation}
The constant $C_j(L)$ depends only on $L$, the dimension, the chosen
integer $a$, the fixed time interval and center cylinder, and smooth
coefficient bounds for the metrics, immersion, Gauss map and cutoffs
on a fixed neighborhood of the boundary compact $K_j$.
It is independent of $\eps,E_0,D_0$, of the interior
compact $K_{j,\eps}'$, and of the auxiliary loss introduced in the proof.
In contrast, the existence radius of the deformations, bounds on their
full coefficients, and bounds on Taylor remainders may depend on
$j,L,E_0,D_0,K_{j,\eps}'$, that loss, the interior margin of the
centers, and the fixed~smooth~test~$\varphi$.
\end{lemma}

\begin{proof}
We give the realization, Taylor justification, and actual mixed-derivative
estimate separately.

\emph{Realization of the first field.}
Let $Z_1,\ldots,Z_b$ be the global holomorphic spanning fields on $Y$.
In the immersion $G$, let $A$ be their column matrix along $f$.
If $r=\dim_\C Y$, enumerate all $r$-row, $r$-column minors $a_i$ of $A$.
On $K_{j,\eps}'$ the sum of their squared moduli has a common positive
lower bound.
Their sup norms and Dirichlet norms are bounded at the fixed parameters:
the minors are smooth holomorphic functions of $X$ on this compact,
and $D(X)$ is bounded.  Set
\[
 \sigma_i=\frac{\overline{a_i}}{\sum_k|a_k|^2},\qquad
 b_{ik}=\sigma_i\bar\partial\sigma_k-
 \sigma_k\bar\partial\sigma_i.
\]
These $(0,1)$ coefficients have bounded $L^2(\Delta)$ norm.  Solve
$\bar\partial Q_{ik}=b_{ik}$, skew-symmetrically, with bounded
$W^{1,2}$ norm.  Explicitly, if $b_{ik}=b_{ik,0}\,d\bar\zeta$,
take $Q_{ik}=4\partial_\zeta\Delta_D^{-1}b_{ik,0}$, with the
convention $\Delta=4\partial_\zeta\partial_{\bar\zeta}$.
Here $\Delta_D^{-1}$ is the zero-Dirichlet inverse on the unit disc.
The unit disc has smooth boundary, the principal coefficients are
the identity, and all lower coefficients vanish.  The scalar
Dirichlet estimate therefore gives
\[
 \|\Delta_D^{-1}b_{ik,0}\|_{W^{2,2}(\Delta)}
 \le C\|b_{ik,0}\|_{L^2(\Delta)}
\]
\cite[Theorem~9.15 and Lemma~9.17]{GilbargTrudinger}, proving the
asserted $W^{1,2}$ bound.  Linearity preserves skew symmetry.
Smooth input on the closed disc gives smooth output there
\cite[Theorem~8.13]{GilbargTrudinger}.

Here is the precise boundary estimate.  For a smooth function $q$
on the closed disc with boundary Fourier coefficients $a_k$, let
$h(re^{i\theta})=\sum_k a_kr^{|k|}e^{ik\theta}$ be its harmonic
extension.  Integration by parts gives
\[
 \int_\Delta|\nabla q|^2
 =\int_\Delta|\nabla h|^2+\int_\Delta|\nabla(q-h)|^2
 \ge2\pi\sum_k|k||a_k|^2.
\]
If $q_0(r)$ is the angular average, then for $1/2\le r\le1$
\[
 |q_0(1)|^2\le2|q_0(r)|^2
          +2(1-r)\int_r^1|q_0'(s)|^2\,ds.
\]
Integrating in $r$ and applying Jensen bounds the constant
coefficient $|a_0|^2$ by $C\|q\|_{W^{1,2}(\Delta)}^2$.
The nonconstant $L^2$ Fourier sum is bounded by its $|k|$-weighted
sum.  Consequently
\[
 \|q|_{\partial\Delta}\|_{L^2}^2+D(q|_{\partial\Delta})^2
 \le C\|q\|_{W^{1,2}(\Delta)}^2.
\]
Applied entrywise to $Q_{ik}$, this bounds its boundary $D$ and
$L^2$ norms.  The operator constants are fixed-disc constants;
the input norms can still depend on $K'_{j,\eps}$.
The functions
\[
 g_i=\sigma_i+\sum_kQ_{ik}a_k
\]
are holomorphic and $\sum_i g_i a_i=1$: differentiating uses
$\sum_k a_k\sigma_k=1$ and
$\sum_k a_k\bar\partial\sigma_k=0$, while skew symmetry cancels the
quadratic term in the sum.

A $W^{1,2}$ trace need not be bounded.  For a large number $P$ let
$\alpha$ be the scalar outer function with positive center value and
boundary modulus
\[
 r_P=(1+|Q|^2/P)^{-1}.
\]
All norms and products involving $Q$ in the following boundary estimates
refer to its trace on the circle.
The maps $Q\mapsto\log(1+|Q|^2/P)$ have Lipschitz constant
$C/\sqrt P$, and their $L^2$ norms are at most $C\|Q\|_2/\sqrt P$.
The maps $Q\mapsto r_PQ$ have bounded Lipschitz constant and sup norm
at most $C\sqrt P$.  The phase of $\alpha$ is the Hilbert transform of
$\log r_P$.  Thus
\[
 \|\alpha\|_\infty\le1,\quad
 D(\alpha)\le C D(Q)/\sqrt P,\quad
 D(\alpha Q)\le C D(Q),\quad
 \|\alpha Q\|_\infty\le C\sqrt P.
\]
In the third estimate the $\sqrt P$ from the bounded product multiplies
the $P^{-1/2}$ bound for the phase.  Moreover,
\[
 \langle1-|\alpha|^2\rangle\le C\|Q\|_2^2/P,
 \qquad -\log\alpha(0)=\langle\log(1+|Q|^2/P)\rangle\longrightarrow0.
\]
Given $\eta>0$, choose $P$ uniformly for the sequence so that these
last two quantities are at most $\eta$ and $D(\alpha)\le1$.
The boundary estimates bound the sup norms and $D$ norms of
$\alpha g_i$.  These functions are holomorphic, so the maximum principle
also bounds their interior sup norms at the fixed parameters; it is
not applied to the nonholomorphic products $\alpha Q$.

For each minor $a_i$, Cramer's rule supplies a vector $c_i$, holomorphic
in the disc, with $Ac_i=a_iUl$.  Namely use the adjugate on its selected
rows and columns and set the other coefficient entries to zero.
The identity follows by expanding the augmented determinants: every
$(r+1)$-minor of $(A,Ul)$ is zero because this matrix still has rank $r$.
This also treats selected minors that vanish identically.
Consequently $c=\sum_i\alpha g_i c_i$ has bounded sup norm and $D$ and
obeys $Ac=\alpha Ul$.  Write $c=(c_1,\ldots,c_b)$ and compose the
complex flows of $Z_1,\ldots,Z_b$
with times $z h c_1,\ldots,z h c_b$, starting at $f$.
For a sufficiently small common complex parameter $z$ this gives a
holomorphic family $f_z$, smooth up to the circle, whose first field is
\begin{equation}
 V=\left.\partial_z f_z\right|_0=\alpha hE_fl.
 \label{var:first-field}
\end{equation}
The existence radius is uniform at the fixed parameters, since the
initial images lie in $K_{j,\eps}'$ and all flow times are bounded.
On the boundary the displacement is $m$ times a uniformly bounded
smooth function.  Because $m\le\delta_i^a$, decreasing this radius
makes each new $\delta_i$ comparable to its old value, with constants
between $1/2$ and $3/2$.  Thus the spatial boundary restrictions persist.
The variable $w$ and all times $t$ are unchanged, so no uniform margin
below the upper time boundary is needed.  The center stays in
$\mathcal C$ by its already established margin.

\emph{Uniform Taylor remainders at the fixed parameters.}
We spell this out because smoothness of individual discs does not by
itself imply the uniform expansion needed for an infimum.
On the controlled interior compact the flow map and all its derivatives
through any fixed finite order are bounded.  Differentiating in the
disc variable introduces one derivative of $f$, $h$, or $c$.
Their Dirichlet norms are bounded; for $f$ use local metric equivalence
on $K_{j,\eps}'$ and \eqref{var:energy}.  Hence the family and its first
three real parameter derivatives have bounded $W^{1,2}$ norms, with
bounded sup norms for their values.  Differentiating the area integrands
through order three yields only bounded coefficients times products of
at most two such first disc derivatives.  Cauchy--Schwarz supplies
uniform integrable bounds.  The same reasoning applies to the Gauss
map, whose derivatives are bounded on the controlled compact.

At the boundary write $\beta=h/m$.  The boundary values of $f_z$ can
be written in ambient coordinates as a smooth function of $X$, $c$,
$\beta$, and $z m$, on a fixed compact range.  More precisely, for
either cutoff,
\[
 \delta_i(f_z)=\delta_i(f)+m R_i(z,X,c,\beta),\qquad R_i(0,\cdot)=0,
\]
where the real parameter derivatives through order three of $R_i$
have bounded sup and $D$ norms at the fixed parameters, by
\eqref{var:boundary-calculus}.  Since
$m/\delta_i=\delta_i^{a-1}\delta_{3-i}^{a}$ is smooth on the closed
boundary region, the function
\[
 \log\delta_i(f_z)-\log\delta_i(f)
 =\log\bigl(1+(m/\delta_i)R_i\bigr)
\]
and its parameter derivatives through order three have uniformly
bounded $D$ norms.  For reciprocal walls, use
\[
 \frac1{\delta_i(f_z)}-\frac1{\delta_i(f)}
 =-\frac{m}{\delta_i^2}\,
 \frac{R_i}{1+(m/\delta_i)R_i}.
\]
The prefactor is smooth and bounded when $a\ge2$.  Therefore all
nonzero parameter derivatives of the wall costs through order three
are uniformly bounded, even when the unchanged wall cost is large.
The remaining boundary integrands $N,d_{0,j}$ have bounded derivatives
on $K_j\times[0,T]$.  Finally the quadratic seminorm of the logarithmic
wall has three bounded parameter derivatives by its Hilbert space
pairing formula.  These observations prove a uniform $O(|z|^3)$
remainder after parameter-circle averaging of the entire functional
minus the test.  The constants in this paragraph are allowed to depend
on the interior compact and on $\eta$.

\emph{The actual mixed energy derivative uses boundary constants only.}
Set $\widehat V=dG\,V=\alpha hUl$ on the circle.  Equations
\eqref{var:boundary-calculus}--\eqref{var:first-field} imply
\begin{equation}
 \|\widehat V\|_\infty\le L,
 \qquad D(\widehat V)^2\le
 C_j(L)\bigl(1+\mathcal E+D(l)^2\bigr).
 \label{var:first-field-bound}
\end{equation}
Here only $\|\alpha\|_\infty\le1$ and $D(\alpha)\le1$ were used;
the corona coefficients have disappeared.

For clarity, consider any one of the closed $(1,1)$ forms $\omega$ in
the three spatial area terms of \eqref{var:energy}.  It is a smooth
closed form on $Y$ after pullback.  In ambient coordinates near the
boundary write its intrinsic tensor as
$i b_{p\bar q}(X)\,dX^p\wedge d\bar X^q$, extending the coefficients
smoothly from $G(\mathcal W_j)$ near $G(K_j)$ as above.  Deformed
boundaries stay intrinsically in $\mathcal W_j$ for a smaller parameter
radius.  Closure is used intrinsically
before this extension.  Cartan's formula and Stokes give
\begin{equation}
 \begin{split}
 \left.\partial_z\partial_{\bar z}
       \int_\Delta f_z^*\omega\right|_0
 =i\int_0^{2\pi}\bigl[&
 b_{p\bar q}(X)\widehat V^p\partial_\theta
                         \overline{\widehat V^q}\\
 &+(\partial_{\bar r}b_{p\bar q})(X)
 \widehat V^p\overline{\widehat V^r}\,
                         \partial_\theta\overline{X^q}\bigr]\,d\theta.
 \end{split}
 \label{var:boundary-area-identity}
\end{equation}
One can also verify the identity using local potentials: the mixed
parameter derivative equals the boundary normal derivative of
$\omega(V,\bar V)$, with the corresponding Stokes normalization.
In \eqref{var:boundary-area-identity}, the holomorphic dependence of
$Gf_z$ eliminates a mixed second parameter field.  The identity is
global; it does not require partitioning $\omega$ into nonclosed forms.

The coefficient $b$ has bounded $C^2$ norm determined only by $K_j$.
For $B=\|\widehat V\|_\infty$,
\[
 \begin{split}
 D(b(X)\widehat V)&\le C_j\bigl(D(\widehat V)+B D(X)\bigr),\\
 D((\bar\partial b)(X)\widehat V\overline{\widehat V})
 &\le C_j\bigl(BD(\widehat V)+B^2D(X)\bigr).
 \end{split}
\]
Pairing each angular derivative by \eqref{var:boundary-calculus}
bounds the absolute value of \eqref{var:boundary-area-identity} by
$C_j(L)(D(X)^2+D(\widehat V)^2)$.  This proves the required estimate
for all spatial areas.  The $w$ area is unchanged.  In particular for
the Euclidean area the identity reduces, up to its fixed normalization,
to $D(\widehat V)^2$.

It remains to check the logarithmic seminorm; its unbounded argument
cannot be handled by an unweighted smooth-composition assertion.
On the boundary put $Z=\alpha\beta Ul$, so that $\widehat V=mZ$.
Then
\[
 \|Z\|_\infty\le L,\qquad
 D(Z)\le C_j(L)(1+\sqrt{\mathcal E}+D(l)).
\]
Regarding $p=\delta_1\delta_2$ as a smooth ambient function, the tensors
\begin{equation}
 \begin{split}
 m\partial\log p&=p^{a-1}\partial p,\\
 m^2\dd\log p&=p^{2a-1}\dd p
                  -p^{2a-2}\partial p\otimes\bar\partial p
 \end{split}
 \label{var:cleared-log-tensors}
\end{equation}
extend smoothly to the closed boundary region.
For $q_z=\log p(f_z)$ their contractions against $Z$ and
$(Z,\bar Z)$ are exactly $\partial_zq_z|_0$ and
$\partial_z\partial_{\bar z}q_z|_0$, respectively.  Consequently each
has $D$ bounded by $C_j(L)(1+\sqrt{\mathcal E}+D(l))$.
Write $b=\partial_zq_z|_0$ and $c=\partial_z\partial_{\bar z}q_z|_0$,
and let $B_D$ be the sesquilinear pairing
$B_D(u,v)=\sum_k|k|u_k\overline{v_k}$.  Since $q_z$ is real,
differentiating this quadratic seminorm gives exactly
\begin{equation}
 \left.\partial_z\partial_{\bar z}D(q_z)^2\right|_0
       =2D(b)^2+2\Re B_D(q,c).
 \label{var:log-energy-identity}
\end{equation}
Since $D(q)^2\le\mathcal E$, Cauchy--Schwarz proves
\begin{equation}
 \left|\left.\partial_z\partial_{\bar z}
                 \mathcal E(f_z,w)\right|_0\right|
 \le C_j(L)(1+\mathcal E+D(l)^2).
 \label{var:actual-energy-derivative}
\end{equation}
All constants in \eqref{var:actual-energy-derivative} use boundary
coefficients only, independently of the realization and Taylor constants.

\emph{Testing the infimum and removing the auxiliary loss.}
Every deformed value of \eqref{var:penalized-functional} is at least
$\mu_{j,\eps}$, while the values at $z=0$ approach it.  Average over
$|z|=r$, subtract the value at zero, and divide by $r^2$.
First take the minimizing-sequence limit with $r$ fixed, and then
$r\downarrow0$.  The uniform remainder just proved yields a
nonnegative lower limit for the mixed derivative of cost minus test.
Its nonpenalty terms are
\[
 \langle|\alpha|^2 T_b(l,\bar l)\rangle
 -|\alpha(0)|^2h(0)^2\varphi_{E,c}(l(0),\overline{l(0)}).
\]
The center frames and their inverses are bounded independently of
interior area by Lemma~\ref{lem:disc-frames}.  Thus the test term is
bounded, as is the boundary symbol by \eqref{var:symbol-bounds}.
Removing $\alpha$ changes these terms by at most $C_{j,L,\varphi}\eta$.
Equation~\eqref{var:actual-energy-derivative} has a coefficient
independent of $\eta$.  Send $\eta\downarrow0$ after the sequence and
radius limits.  This proves \eqref{var:deformation-estimate} with the
stated uniformities.  In particular, none of the constants controlling
the shrinking realization radius multiplies the final $\eps\mathcal E$
error.
\end{proof}

\subsection{Removing the negative part and factoring the boundary form}

The allowance for disc-dependent columns in
Lemma~\ref{lem:disc-deformation} is essential here.  Let
$S=(T_b)_-$ be the positive semidefinite negative part of $T_b$.
Matrix clipping is Lipschitz in the Frobenius norm, so
$\|S\|_\infty\le C_j$ and $D(S)^2\le C_j(1+\mathcal E)$.
Let $S_P$ be its Fej\'er mean of order $P$.  Positivity of the Fej\'er
kernel preserves its sup bound, and its Fourier multipliers give
\begin{equation}
 \|S_P-S\|_2^2\le C_j(1+\mathcal E)/P.
 \label{var:fejer-error}
\end{equation}
Each column $l_i(\zeta)=\zeta^{P+1}S_P(\zeta)e_i$ is holomorphic,
vanishes at zero, and has
\[
 \|l_i\|_\infty\le C_j,\qquad D(l_i)^2\le C_j(P+1).
\]
The frequency shift has modulus one on the circle, so summing
\eqref{var:deformation-estimate} over these columns tests
$\tr(T_bS_P^2)$.  In comparison,
$\tr(T_bS^2)=-\tr(S^3)$.
For fixed $j,\eps$, discard finitely many sequence terms and let
$B_\eps$ bound $1+\mathcal E_\nu$, with
$\eps B_\eps\to0$ as $\eps\downarrow0$, as permitted by
\eqref{var:vanishing-penalty}.  Choose an integer $P=P_\eps$, fixed
along this sequence, with
\[
 B_\eps/P_\eps\longrightarrow0,
 \qquad \eps P_\eps\longrightarrow0;
\]
for example round $\sqrt{B_\eps/\eps}$ upward.
The trace error is at most $C_j\|S_P-S\|_2$ by boundedness of both
matrices.  Equations~\eqref{var:deformation-estimate} and
\eqref{var:fejer-error} therefore imply
\begin{equation}
 \lim_{\eps\downarrow0}\limsup_\nu
       \langle\tr((T_b)_-^3)\rangle=0,
 \qquad
 \lim_{\eps\downarrow0}\limsup_\nu\|(T_b)_-\|_1=0.
 \label{var:negative-part}
\end{equation}
The second conclusion follows by H\"older's inequality in fixed dimension.
This argument uses both sides of the intermediate frequency choice;
a bound merely of the form $\eps\mathcal E=O(1)$ would not suffice.

Fix $\delta>0$.  Apply the smooth scalar function
$\lambda\mapsto\delta+(\lambda+\sqrt{\lambda^2+\delta^2})/2$
to $T_b$ and call the resulting positive boundary matrix $L$.
Then
\[
 L\ge T_b+\delta I,\qquad L\ge\delta I,
 \qquad \|L-(T_b)_+\|\le\tfrac32\delta.
\]
Functional calculus in fixed dimension gives
$D(L)\le C_{j,\delta}D(T_b)$.
Lemma~\ref{lem:boundary-factorization} gives
$L=H^*H$, with $H,H^{-1}$ holomorphic and smooth on the closed disc;
this is the smooth boundary factorization of
\cite[Theorem~2.1]{BerndtssonRosay}, with the quantitative seminorm
estimate established in Section~\ref{sec:discs}.
For every constant unit column $e$,
\[
 l=H^{-1}e,\qquad \|l\|_\infty\le\delta^{-1/2},
 \qquad D(l)^2\le C_{j,\delta}(1+\mathcal E).
\]
These are permissible disc-dependent tests at fixed $j,\eps,\delta$.
Since $T_b\le L$, their boundary quadratic costs are at most one.
Equation~\eqref{var:vanishing-penalty} now gives
\begin{equation}
 h(0)^2\varphi_{E,c}\le(1+o(1))H(0)^*H(0)
 \label{var:center-majorant}
\end{equation}
as the sequence limit and then $\eps\downarrow0$ are taken at fixed
$j,\delta$.  To see that this is a matrix assertion, if the largest
violation did not tend to zero, choose its unit eigenvector separately
for each disc.  Those columns satisfy exactly the same sup and $D$
bounds, contradicting Lemma~\ref{lem:disc-deformation}.
Thus no fixed-test assumption has entered either the Fej\'er test or
the spectral-factor test.

\subsection{The backward clock and the symplectic cost}

Choose cutoffs $0\le\chi_i\le1$ which are one wherever $\delta_i$
differs, or starts to differ, from one.  Require
\[
 \supp\chi_1\subset\{u>A_j/3\},\qquad
 \supp\chi_2\subset\{N(y,0)>c_2B_j^2\}
\]
for a fixed $c_2>0$.  Before choosing $B_j$ we can choose $C'_j\ge1$
depending on the base compact, $d_{0,j}$ and the time interval, and a
fixed integer $c_3$, such that
\begin{equation}
 \begin{split}
 \bigl|p(W_b)-\partial_tN
   -\Delta_{h_t}\bigl(d_{0,j}+(1+s)k_1\bigr)\bigr|
 \le C'_j(1+B_j)^{c_3}
            (\delta_1\delta_2)^{-c_3}\boldone_{\{\chi_2=1\}}.
 \end{split}
 \label{var:fiber-wall-error}
\end{equation}
Indeed Lemma~\ref{lem:symplectic-evolution} gives the exact equality
without the second wall.  The remaining expression is a polynomial
of degree at most two in the Hessian entries.  On a compact base,
$N(y,t)$ and $N(y,0)$ are quadratic in the fiber variables, and their
fixed finite derivatives grow polynomially in those variables.
Differentiating the reciprocal cutoff twice introduces at most
three inverse powers of $\delta_2$; the analogous first-wall factors
have at most three inverse powers of $\delta_1$.
Enlarging a fixed integer $c_3$ bounds all these finitely many terms.
This proves \eqref{var:fiber-wall-error}, including that $C'_j$ is
chosen before $B_j$.
Also
\begin{equation}
 |\Delta_{h_t}k_1|\le C e^\rho\delta_1^{-3}
                         \boldone_{\{\chi_1=1\}}.
 \label{var:base-wall-error}
\end{equation}
Here the controlled gradient and complex Hessian of $u$ bound
$|\dd k_1|_g$ by $C\delta_1^{-3}$ uniformly in large $A_j$.
Since $0<h_t\le g$, each inverse eigenvalue relative to $g$ is at
most $e^\rho$, proving \eqref{var:base-wall-error}.

Fix integers $a,c$ sufficiently large that $4a\ge c_3+3$ and
$c>c_3+3$, retaining all earlier requirements on $a$.  Put
\begin{equation}
 \begin{split}
 \gamma_0={}&m^4\exp\bigl[-\chi_1(\rho+c\log(j+2))\bigr]\\
 &\mathrel{}\cdot\exp\bigl[-\chi_2(
       \log C'_j+c\log(j+2)+c\log(1+B_j))\bigr],\\
 \gamma={}&\gamma_0/\langle\gamma_0\rangle.
 \end{split}
 \label{var:clock-weight}
\end{equation}
Choose $B_j$ so large that $\log C'_j/B_j^2\to0$.
The preceding choice of $C'_j$ makes this noncircular.
Since $-\log\delta_i\le k_i$, \eqref{var:small-walls} controls
$\langle-\log m\rangle$.  On $\supp\chi_1$,
\eqref{var:tail-wall} controls both $\rho$ and $j$ by $d_{0,j}$.
On $\supp\chi_2$, the pointwise boundary inequality
$N(f,t)\ge N(f,0)$ and \eqref{var:small-walls} give
$\langle\chi_2\rangle\le C/B_j^2$.
It follows, in the stated successive limits, that
\begin{equation}
 \langle-\log\gamma_0\rangle\longrightarrow0,
 \qquad \langle\gamma_0\rangle\longrightarrow1,
 \qquad \langle\log\gamma\rangle\longrightarrow0.
 \label{var:clock-loss}
\end{equation}
For example, the second limit also follows from
$e^{\langle\log\gamma_0\rangle}\le\langle\gamma_0\rangle\le1$.

The powers and exponentials in \eqref{var:clock-weight} bound the
weighted errors in \eqref{var:fiber-wall-error} and
\eqref{var:base-wall-error} by a quantity tending to zero.  On the
support of each error its corresponding cutoff equals one, so the
cancellation is literal; $m^4$ clears the displayed inverse wall
powers, $e^{-\rho}$ clears $e^\rho$, and
$(C'_j)^{-1}(1+B_j)^{-c}(j+2)^{-c}$ clears the fiber error.
The denominator $\langle\gamma_0\rangle$ tends to one, and $t/s$ is
bounded because centers stay in $\mathcal C$.
Since $(\partial_t-\Delta_{h_t})d_{0,j}\ge0$, we obtain
\begin{equation}
 \left\langle\frac ts\gamma(\partial_tN+\partial_t d_{0,j})
                      +k_1+k_2\right\rangle
 \ge\left\langle\frac ts\gamma p(W_b)\right\rangle-o_j(1).
 \label{var:clock-comparison}
\end{equation}
Here $o_j(1)$ tends to zero after the earlier limits.

The clock test giving the left side of \eqref{var:clock-comparison}
must run backward.  At fixed $j$,
\[
 \|\gamma\|_\infty\le C_j,\qquad
 D(\gamma)^2\le C_j(1+\mathcal E).
\]
These follow by \eqref{var:boundary-calculus} from the smooth
coefficients in \eqref{var:clock-weight}; the denominator has a
positive lower bound by Jensen and the bounded cost.
For an explicit bound, let $M$ bound the unpenalized boundary cost
along the fixed-$j$ sequences.  Then
$\langle k_1+k_2\rangle\le M$,
$\langle-\log p\rangle\le M$, and
$\langle\chi_1\rho\rangle\le M$ by the tail wall.
Hence
\[
 \langle-\log\gamma_0\rangle
 \le (4a+1)M+2c\log(j+2)+\log C'_j+c\log(1+B_j),
\]
and Jensen gives a positive lower bound for
$\langle\gamma_0\rangle$ depending only on these fixed-$j$ data.
No lower bound on the individual cutoffs or boundary times is used.
Let $\Gamma$ be holomorphic with boundary real part $\gamma-1$ and
$\Gamma(0)=0$.  For $s'<s$ set
\[
 w'=w+\tfrac12\log(s'/s)\Gamma.
\]
On the circle the new time is $t(s'/s)^\gamma$.
In the interior its exponent is the positive harmonic extension of
$\gamma$.  Hence every time decreases, including when the original
disc has arbitrarily little margin below $T$.
The center of $w'$ is still zero.

Comparing the perturbed functional with its infimum, first along the
minimizing sequence and then letting $s'\uparrow s$, gives
\begin{equation}
 \partial_s\varphi(f(0),s)
 \ge\left\langle\frac ts\gamma(\partial_tN+\partial_t d_{0,j})
                       +k_1+k_2\right\rangle
       -C_j\eps(1+\mathcal E)-o(1).
 \label{var:backward-test}
\end{equation}
The sign follows because the quotient denominator $s'-s$ is negative.
For the energy error only the $w$ area changes, and
$D(\Gamma)\le C D(\gamma)$ gives its derivative bound by
Cauchy--Schwarz.  Its Taylor expansion is quadratic in $\log(s'/s)$.
For the boundary cost, derivatives of $t(s'/s)^\gamma$ are uniformly
bounded at fixed $j$; $N,d_{0,j}$ are smooth down to $t=0$ on $K_j$.
The possibly unbounded unchanged walls multiply $1+s'$ affinely and
thus have no second time remainder.  These facts justify the
sequence-then-difference limit in \eqref{var:backward-test} without a
bound on the imaginary part of $\Gamma$.

Let $\mathcal J_E$ be the bivector components in $E_f$.  Its holomorphic
sup bound depends only on $K_j$ by Lemma~\ref{lem:disc-frames}.
Since $p$ is quadratic in the Hermitian form,
$p(W_b)=m^{-4}p(T_b,\mathcal J_E)$.
At fixed $j$, the factor $\gamma m^{-4}$ is bounded.  Therefore
\eqref{var:negative-part}, the uniform symbol bounds, and
$L=(T_b)_++O(\delta)$ allow the replacement
\begin{equation}
 \left\langle\frac ts\gamma p(W_b)\right\rangle
 =\left\langle\frac ts\gamma m^{-4}p(L,\mathcal J_E)\right\rangle
       +O_j(\delta)+o(1)
 \label{var:positive-replacement}
\end{equation}
after $\eps\downarrow0$.
The vector $(\bigwedge^2H)\mathcal J_E$ is holomorphic and nonzero.
Logarithmic submean for its squared norm, followed by scalar Jensen,
gives
\begin{equation}
 \begin{split}
 \left\langle\frac ts\gamma m^{-4}p(L,\mathcal J_E)\right\rangle
 &\ge \exp\langle\log\gamma\rangle\,
 h(0)^{-4}p(H(0)^*H(0),\mathcal J_E(0)).
 \end{split}
 \label{var:symplectic-submean}
\end{equation}
Indeed $\langle\log(t/s)\rangle=2\Re w(0)=0$ and
$\log h(0)=\langle\log m\rangle$, which give precisely the factors
in \eqref{var:symplectic-submean}.  If desired, logarithmic submean
follows by applying submean to $\log(|Z|^2+\eta)$ and then
$\eta\downarrow0$ for the displayed holomorphic vector $Z$.

\subsection{Passing from positive majorants to the required form}

\begin{lemma}[Symplectic cost and the Schur complement]
\label{lem:symplectic-schur}
In a cotangent vector space with canonical bivector
$\mathcal J=\sum_i e_i\wedge f_i$, let $\mathcal H>0$ be Hermitian,
and suppose $\mathcal H\ge W$, where the vertical block $C$ of $W$
is positive.  Then
\begin{equation}
 S=\Schur(\mathcal H)\ge0,\qquad
 S\ge\Schur(W),\qquad p(\mathcal H)\ge\tr(C^{\mathsf T}S).
 \label{var:schur-cost}
\end{equation}
\end{lemma}

\begin{proof}
Write $\mathcal H=\left(\begin{smallmatrix}A&B\\B^*&D\end{smallmatrix}\right)$,
$S=A-BD^{-1}B^*$ and $R=D^{-1}B^*$.
The complex-linear change $(x,y)\mapsto(x,y+Rx)$ makes the form
block diagonal with blocks $(S,D)$.  It carries the canonical bivector
to
\[
 \sum_i e_i\wedge f_i+\sum_{k,i}R_{ki}f_k\wedge f_i.
\]
Its cross and vertical components are orthogonal for the induced
exterior-square form, whence
\[
 p(\mathcal H)=\tr(SD^{\mathsf T})+
       \left|\sum_{k,i}R_{ki}f_k\wedge f_i\right|_D^2.
\]
The shear need not be symplectic; its additional vertical term is
retained and is nonnegative.  Since $D\ge C$ and $S>0$, the last
identity gives $p(\mathcal H)\ge\tr(SC^{\mathsf T})$.
Finally, for each horizontal $x$,
\[
 S(x,\bar x)=\inf_y\mathcal H((x,y),\overline{(x,y)})
 \ge\inf_y W((x,y),\overline{(x,y)})
 =\Schur(W)(x,\bar x).
\]
The last infimum is finite and has the usual Schur expression because
$C>0$.  This proves the Lemma.
\end{proof}

\begin{proof}[Completion of the proof of Proposition~\ref{prop:disc-inequality}]
Combine \eqref{var:center-majorant}, \eqref{var:clock-comparison},
\eqref{var:backward-test}, \eqref{var:positive-replacement}, and
\eqref{var:symplectic-submean}.  More explicitly, if $r\ge0$ denotes
the error tending to zero in \eqref{var:center-majorant}, the form
with frame matrix
\[
 \mathcal H_E=(1+r)h(0)^{-2}H(0)^*H(0)
\]
is a positive majorant of the test Hessian at its center, after an
arbitrarily small enlargement if necessary.  Transform it to the
fixed cotangent coordinates.  Its symplectic cost is
\[
 p(\mathcal H)=(1+r)^2h(0)^{-4}
                 p(H(0)^*H(0),\mathcal J_E(0)).
\]
For fixed $j,\delta$ first take the minimizing-sequence limits and
$\eps\downarrow0$; then send $\delta\downarrow0$; finally let
$j\to\infty$.  Equations~\eqref{var:small-walls} and
\eqref{var:clock-loss} give a diagonal sequence of approaching centers
and positive majorants $\mathcal H$ such that
\[
 \limsup p(\mathcal H)\le\partial_s\varphi(y_*,s_*).
\]
No compactness of the complete matrices $\mathcal H$ is needed.
Their test vertical blocks $C_c$ converge to $C>0$, so for late
centers $C_c\ge cI$ with a fixed $c>0$.
Lemma~\ref{lem:symplectic-schur} bounds their positive Schur forms
$S_c$ by $c\tr S_c\le p(\mathcal H)$.
Pass to a convergent subsequence $S_c\to K$.  The same Lemma and
continuity of the test Hessian and its Schur complement give
\[
 K\ge0,\qquad K\ge\Schur(W),\qquad
 \tr(C^{\mathsf T}K)\le\partial_s\varphi(y_*,s_*).
\]
This is \eqref{var:viscosity-inequality} and proves the Proposition.
\end{proof}
